\heading{量子力学A-24秋-期中1}
\noteinfo{原资料为手写期中题解；整理时保留题目与解答，按 $m=\hbar\equiv 1$ 处理。}

\begin{question}
本题为若干基础概念与 Ehrenfest 定理的简答题。
\begin{enumerate}
  \item 在动量表象下，写出 $\{\hat x,\hat y,\hat p_z\}$ 的共同本征函数，其对应本征值为 $\{x_0,y_0,p_0\}$。
  \item 对势函数 $V(\vb r)=V_0\theta(z)$，列出体系的守恒量，并说明其对应的对称性。
  \item 一维粒子的初始位置、动量平均值分别为 $x_0,p_0$。分别对下列两个哈密顿量，求 $\bar x(t)$ 与 $\bar p(t)$：
  \begin{enumerate}
    \item $H={\hat p^2\over 2}-\hat x$；
    \item $H={\hat p^2\over 2}-\hat p$。
  \end{enumerate}
\end{enumerate}
\begin{solution}
\begin{enumerate}
  \item 在动量表象中 $\hat x=\ii\partial_{p_x}$、$\hat y=\ii\partial_{p_y}$，故
  \[
    \Phi_{x_0,y_0,p_0}(p_x,p_y,p_z)=C\exp[-\ii(x_0p_x+y_0p_y)]\delta(p_z-p_0),
  \]
  归一化常数按 $\delta$ 归一化约定选取。
  \item 势能只依赖 $z$，所以平移对称给出 $p_x,p_y$ 守恒；绕 $z$ 轴转动对称给出 $L_z$ 守恒；若势不显含时间，则 $H$ 守恒。
  \item 用 Ehrenfest 定理。
  \begin{enumerate}
    \item 对 $H={\hat p^2/2}-\hat x$，有 $\dot{\bar p}=1$、$\dot{\bar x}=\bar p$，故
    \[
      \bar p(t)=p_0+t,
      \qquad \bar x(t)=x_0+p_0t+{t^2\over 2}.
    \]
    \item 对 $H={\hat p^2/2}-\hat p$，有 $\dot{\bar p}=0$、$\dot{\bar x}=\bar p-1$，故
    \[
      \bar p(t)=p_0,
      \qquad \bar x(t)=x_0+(p_0-1)t.
    \]
  \end{enumerate}
\end{enumerate}
\end{solution}
\end{question}

\begin{question}
一维束缚态势函数为 $V(x)=4x^4$。设体系处于基态。
\begin{enumerate}
  \item 判断基态宇称；
  \item 求坐标、动量均值；
  \item 用不确定关系估计基态能量。
\end{enumerate}
\begin{solution}
势能为偶函数，非简并基态可取为偶宇称。
因此
\[
  \expval{x}=0,
  \qquad \expval{p}=0 .
\]
令 $\Delta x=s$，用 $\Delta x\Delta p\simeq 1/2$ 估计
\[
  E(s)\simeq {1\over 2}(\Delta p)^2+4(\Delta x)^4
  ={1\over 8s^2}+4s^4.
\]
极小条件 $E'(s)=0$ 给出 $s=1/2$，故
\[
  E_0\simeq {3\over 4}.
\]
\end{solution}
\end{question}

\begin{question}
无限深方势阱的范围为 $0<x<2a$，阱内势能为零、阱外势能为无穷大。某时刻粒子的波函数为
\[
  \psi(x)=A\sin {\pi x\over a}\,\theta(a-x).
\]
\begin{enumerate}
  \item 求 $A$；
  \item 求能量均值；
  \item 展开到宽度为 $2a$ 的势阱本征态后，哪些能量本征值有非零测量概率？
  \item 哪个能量本征值的测量概率最大？最大概率为多少？
\end{enumerate}
\begin{solution}
\begin{enumerate}
  \item 归一化给出
  \[
    1=A^2\int_0^a\sin^2{\pi x\over a}\dd{x}=A^2{a\over 2},
    \qquad A=\sqrt{2\over a}.
  \]
  \item 在势阱内 $H=-\frac12\dv[2]{}{x}$，因此
  \[
    \expval{H}= {1\over 2}\int_0^a |\psi'(x)|^2\dd{x}
    ={\pi^2\over 2a^2}.
  \]
  \item 宽度为 $2a$ 的本征态和能级为
  \[
    \phi_n(x)={1\over \sqrt a}\sin {n\pi x\over 2a},
    \qquad E_n={n^2\pi^2\over 8a^2},\quad n=1,2,\ldots .
  \]
  展开系数
  \[
    c_n=\sqrt{2\over a^2}\int_0^a\sin {\pi x\over a}\sin {n\pi x\over 2a}\dd{x}.
  \]
  因而 $n=2$ 必可出现；奇数 $n$ 也可出现；除 $n=2$ 外的偶数项为零。更具体地，
  \[
    c_2={1\over\sqrt2},\qquad
    c_n={4\sqrt2\over \pi}{\sin(n\pi/2)\over 4-n^2}\quad(n\ne 2).
  \]
  \item 最大概率对应 $n=2$，
  \[
    E_2={\pi^2\over 2a^2},\qquad P_2=|c_2|^2={1\over 2}.
  \]
\end{enumerate}
\end{solution}
\end{question}

\begin{question}
一维散射势为
\[
  V(x)=\alpha\delta(x)+V_0\theta(x),\qquad \alpha<0,
\]
能量 $E>V_0$ 的自由粒子从左方入射。求透射系数 $T$。
\begin{solution}
记
\[
  k=\sqrt{2E},\qquad K=\sqrt{2(E-V_0)} .
\]
取
\[
  \psi(x)=\begin{cases}
  \ee^{\ii kx}+r\ee^{-\ii kx},&x<0,\\
  t\ee^{\ii Kx},&x>0.
  \end{cases}
\]
边界条件为 $1+r=t$ 及
\[
  \psi'(0^+)-\psi'(0^-)=2\alpha\psi(0).
\]
解得
\[
  t={2\ii k\over \ii(k+K)-2\alpha}.
\]
故透射系数为透射流与入射流之比：
\[
  T={K\over k}|t|^2={4kK\over (k+K)^2+4\alpha^2}.
\]
\end{solution}
\end{question}

\begin{question}
一维谐振子的哈密顿量为 $\hat H=(\hat p^2+\hat x^2)/2$。相干态满足 $\hat a\ket c=c\ket c$。定义算符
\[
  \hat O=\hat x\hat p+\hat p\hat x.
\]
\begin{enumerate}
  \item 求谐振子第二激发态 $\ket2$ 中 $O$ 的平均值；
  \item 求相干态 $\ket c$ 中 $O$ 的平均值；
  \item 求粒子数算符 $N=a^\dagger a$ 的 $\bar N$、$\overline{N^2}$ 和 $\Delta N$。
\end{enumerate}
\begin{solution}
由
\[
  \hat x={\hat a+\hat a^\dagger\over \sqrt2},\qquad
  \hat p={\hat a-\hat a^\dagger\over \ii\sqrt2}
\]
得
\[
  \hat O=\ii\left[(\hat a^\dagger)^2-\hat a^2\right].
\]
故对数态 $\ket n$，$\mel{n}{O}{n}=0$，特别地第二激发态 $\ket2$ 下 $\bar O=0$。
相干态中
\[
  \expval{O}_c=\ii\left[(c^*)^2-c^2\right].
\]
若 $c=u+\ii v$，则 $\expval{O}_c=4uv$。
粒子数为泊松分布，故
\[
  \bar N=|c|^2,
  \qquad \overline{N^2}=|c|^4+|c|^2,
  \qquad \Delta N=|c|.
\]
\end{solution}
\end{question}
